\section{El producto de números naturales}

Al igual que la adición, el producto es una operación binaria: cada vez que se aplica recibe dos números.

Sobre los números naturales podemos representarlo mediante
\[
\cdot\colon\mathbb{N}\times\mathbb{N}\to\mathbb{N}
\]
Dados dos números naturales \(a\) y \(b\), escribimos su producto como
\[
a\cdot b
\]
Los números \(a\) y \(b\) reciben el nombre de \emph{factores}.
Cuando \(a\) es un número natural, la expresión
\[
a\cdot b
\]
puede interpretarse como una suma en la que \(b\) aparece \(a\) veces:
\[
a\cdot b
=
\underbrace{b+b+\cdots+b}_{a\text{ veces}}
\]
Por ejemplo,
\[
3\cdot5
=
5+5+5
=
15
\]
mientras que
\[
4\cdot2
=
2+2+2+2
=
8
\]

Esta interpretación nos permite comprender inicialmente el producto utilizando una operación que ya conocemos: la adición.

\begin{note}
La interpretación del producto como suma repetida resulta especialmente útil cuando trabajamos con números naturales.

Más adelante ampliaremos el producto a otros conjuntos numéricos. Allí encontraremos expresiones como
\[
\frac12\cdot\frac34
\]
para las cuales ya no tendría sentido hablar literalmente de ``sumar una cantidad media vez''.

Por esta razón, la suma repetida debe entenderse como una forma de construir y comprender inicialmente el producto sobre los números naturales, no como una descripción suficiente de todos los productos que estudiaremos.
\end{note}

\subsection{Producto y factores}

En una suma como
\[
3+5
\]
los números \(3\) y \(5\) reciben el nombre de \emph{sumandos}.
En un producto como
\[
3\cdot5
\]
los números \(3\) y \(5\) reciben el nombre de \emph{factores}.

El valor de la expresión completa recibe el nombre de \emph{producto}. Así,
\[
3\cdot5=15
\]

De la misma manera que una suma puede aparecer como parte de una expresión mayor, un producto también puede constituir solamente una parte de ella.
Por ejemplo, en
\[
7+3\cdot5
\]
tenemos una suma formada por dos términos:
\[
7
\qquad\text{y}\qquad
3\cdot5
\]
Dentro del segundo término aparece, a su vez, un producto cuyos factores son
\[
3
\qquad\text{y}\qquad
5.
\]
Por lo tanto, conviene distinguir el nivel de la expresión que estamos observando:
\[
\underbrace{7}_{\text{sumando}}
+
\underbrace{\left(
  \underbrace{3}_{\text{factor}}
  \cdot
  \underbrace{5}_{\text{factor}}
\right)}_{\text{sumando}}
\]
El producto completo \(3\cdot5\) es uno de los sumandos de la suma exterior. Los números \(3\) y \(5\), en cambio, son factores de ese producto.

\subsection{Una convención para evitar demasiados paréntesis}

Si quisiéramos mostrar explícitamente la estructura de
\[
7+3\cdot5
\]
podríamos escribir
\[
7+(3\cdot5)
\]
Sin embargo, utilizar paréntesis continuamente haría muchas expresiones innecesariamente difíciles de leer.

Por convenio, cuando aparecen productos y sumas en una misma expresión, los productos se consideran agrupados antes que las sumas. Así,
\[
7+3\cdot5
\]
se interpreta como
\[
7+(3\cdot5)
\]
y no como
\[
(7+3)\cdot5
\]
Esta convención recibe el nombre de \emph{precedencia} o \emph{jerarquía de las operaciones}.

No significa que el producto sea una operación ``más importante'' que la suma. Es simplemente una regla de escritura que permite reconocer sin ambigüedad cuáles son los términos de una expresión.

Por ejemplo, en
\[
2+4\cdot3+5
\]
los términos de la suma son
\[
2,\qquad4\cdot3,\qquad5
\]
Si evaluamos el producto,
\[
4\cdot3=12
\]
la expresión queda
\[
2+12+5
\]
Ahora resulta evidente que estamos ante una suma cuyos sumandos son
\[
2,\qquad12,\qquad5
\]
Por lo tanto,
\[
2+4\cdot3+5
=
2+12+5
=
19
\]
